%% ----------------------------------------------------------------------------
%%  §9.2  Primitive. Formula fondamentale del calcolo integrale
%% ----------------------------------------------------------------------------
\section{Primitive. Formula fondamentale del calcolo integrale}
\label{sec:primitive}

\begin{definizione}[Primitiva]\label{def:primitiva}\index{primitiva}
Sia $f$ una funzione definita in un intervallo $I$. Una funzione $F$ si dice una
\emph{primitiva} di $f$ in $I$ se $F$ è derivabile in $I$ e $F'(x)=f(x)$ per ogni $x\in I$.
\end{definizione}

Ad esempio, una primitiva di $f(x)=x$ in $\R$ è $F(x)=\frac{x^{2}}2$; una primitiva di
$\sen x$ è $-\cos x$; una primitiva di $\frac1x$ nell'intervallo $(0,+\infty)$ è $\log x$.
Con la definizione di primitiva, il teorema fondamentale si enuncia così: \emph{se $f$ è
continua in $[a,b]$, la funzione integrale $F(x)=\int_{a}^{x}f(t)\,dt$ è una primitiva di
$f$}. Più in generale, se $f$ è continua in un intervallo $I$ qualsiasi (aperto,
semiaperto o illimitato) e $c\in I$, la funzione $\int_{c}^{x}f(t)\,dt$ è una primitiva di
$f$ in $I$: basta applicare il \cref{teo:fondamentale} in un intervallo chiuso e
limitato contenuto in $I$ che contenga $c$ e il punto $x$ considerato, con $x$ al suo
interno se $x$ non è un estremo di $I$. Dunque:

\begin{center}
\emph{ogni funzione continua in un intervallo ammette primitive.}
\end{center}

Non tutte le funzioni, però, ammettono primitive, neppure se sono integrabili. La
funzione $f$ che vale $-1$ per $x<0$ e $1$ per $x\ge0$ è integrabile in $[-1,1]$, ma non ha
primitive: se fosse $G'=f$ in $[-1,1]$, la derivata $G'$ non si annullerebbe mai e
cambierebbe segno, contro il lemma di Darboux (\cref{lem:darboux}). Viceversa, esistono
funzioni discontinue che ammettono primitive, come la derivata di $x^{2}\sen\frac1x$
(\cref{oss:f2-non-C1}).

\subsection{Le primitive di una funzione in un intervallo}

Se $F$ è una primitiva di $f$, anche $F+c$ lo è, qualunque sia la costante $c$, perché la
derivata di una costante è nulla. Il teorema seguente afferma che in un intervallo non
ce ne sono altre.

\begin{teorema}[Caratterizzazione delle primitive in un intervallo]
\label{teo:caratterizzazione-primitive}
Se $F$ e $G$ sono due primitive di una stessa funzione $f$ in un intervallo $I$, esiste
una costante $c$ tale che
\begin{equation}\label{eq:primitive-costante}
  G(x)=F(x)+c \qquad\perogni x\in I .
\end{equation}
\end{teorema}

\begin{schema}
\Ipotesi{$I$ intervallo; \ $F'=G'=f$ in $I$.}
\Tesi{$G-F$ è costante in $I$.}
\end{schema}

\begin{strategia}
La differenza $H=G-F$ ha derivata nulla in $I$. Una funzione con derivata nulla in un
intervallo è costante: è una conseguenza del teorema di Lagrange, già dimostrata nel
\cref{cap:applicazioni-derivate}. L'ipotesi che $I$ sia un intervallo serve proprio a
poter applicare il teorema di Lagrange tra due punti qualsiasi.
\end{strategia}

\begin{dimostrazione}
Sia $H=G-F$; allora $H$ è derivabile in $I$ e $H'(x)=G'(x)-F'(x)=f(x)-f(x)=0$ per ogni
$x\in I$. Fissato $x_{0}\in I$, sia $x\in I$ con $x\neq x_{0}$. Poiché $I$ è un intervallo,
contiene l'intervallo chiuso di estremi $x_{0}$ e $x$, nel quale $H$ è continua (perché
derivabile) con derivata nulla; per il \cref{cor:funzioni-costanti}, $H$ è costante in
tale intervallo, e quindi $H(x)=H(x_{0})$. Posto $c=H(x_{0})$, si ha $G(x)=F(x)+c$ per ogni
$x\in I$.
\end{dimostrazione}

L'ipotesi che il dominio sia un intervallo è essenziale. Le funzioni $F(x)=\log\abs x$ e
\[
  G(x)=\begin{cases} \log\abs x+1 & \text{se } x>0,\\ \log\abs x+2 & \text{se } x<0 \end{cases}
\]
hanno entrambe derivata $\frac1x$ in $\R\setminus\set{0}$ (lo verificheremo nella
\eqref{eq:derivata-log-abs}), ma la loro differenza non è costante: su ciascuno dei due
intervalli $(-\infty,0)$ e $(0,+\infty)$ la costante può essere scelta indipendentemente.

\begin{figure}[htb]
\centering
\begin{tikzpicture}
\begin{axis}[name=p1, width=0.52\textwidth, height=5cm, axis lines=middle,
  xmin=-0.4, xmax=6.6, ymin=-2.3, ymax=3.4, xtick=\empty,
  ytick=\empty, title={\scriptsize primitive di $\cos x$: $\sen x+c$}, clip=false]
  \addplot[curva3, domain=-0.3:6.4, samples=100] {sin(deg(x))-1};
  \addplot[curva, domain=-0.3:6.4, samples=100] {sin(deg(x))};
  \addplot[curva2, domain=-0.3:6.4, samples=100] {sin(deg(x))+1};
  \addplot[curva3, domain=-0.3:6.4, samples=100] {sin(deg(x))+2};
  \addplot[tratteggio] coordinates {(1.2,-2.2) (1.2,3.3)};
  \addplot[grigio, thick] coordinates {(0.4,-0.358) (2.0,0.222)};
  \addplot[grigio, thick] coordinates {(0.4,0.642) (2.0,1.222)};
  \addplot[grigio, thick] coordinates {(0.4,1.642) (2.0,2.222)};
  \addplot[grigio, thick] coordinates {(0.4,2.642) (2.0,3.222)};
  \addplot[only marks, mark=*, mark size=1.1pt, grigio] coordinates
    {(1.2,-0.068) (1.2,0.932) (1.2,1.932) (1.2,2.932)};
  \node[font=\scriptsize, anchor=north] at (axis cs:1.2,-2.25) {$x_{0}$};
\end{axis}
\begin{axis}[at={(p1.east)}, anchor=west, xshift=0.7cm, width=0.46\textwidth,
  height=5cm, axis lines=middle, xmin=-0.3, xmax=3.6, ymin=-0.2, ymax=1.35,
  xtick={1.5708,3.1416}, xticklabels={$\frac\pi2$,$\pi$}, ytick={1},
  title={\scriptsize $\int_{0}^{\pi}\sen x\,dx=2$}]
  \addplot[draw=none, fill=blu!22, domain=0:3.1416, samples=60] {sin(deg(x))}
    \closedcycle;
  \addplot[curva, domain=0:3.1416, samples=80] {sin(deg(x))};
  \node[font=\scriptsize, blu] at (axis cs:1.5708,0.42) {area $=2$};
\end{axis}
\end{tikzpicture}
\caption{A sinistra, alcune primitive di $\cos x$: i loro grafici si ottengono l'uno
dall'altro per traslazione verticale, e nello stesso punto $x_{0}$ hanno tangenti
parallele, di pendenza $\cos x_{0}$. A destra, l'area sotto un arco di sinusoide, calcolata
con la formula fondamentale \eqref{eq:formula-fondamentale}.}
\label{fig:primitive}
\end{figure}

\subsection{La formula fondamentale}

Il risultato seguente riconduce il calcolo degli integrali definiti alla ricerca delle
primitive.

\begin{teorema}[Formula fondamentale del calcolo integrale]\label{teo:formula-fondamentale}
\index{formula fondamentale del calcolo integrale}
Sia $f$ una funzione continua in $[a,b]$ e sia $G$ una primitiva di $f$ in $[a,b]$. Allora
\begin{equation}\label{eq:formula-fondamentale}
  \int_{a}^{b}f(x)\,dx=\Bigl[G(x)\Bigr]_{a}^{b}=G(b)-G(a) .
\end{equation}
\end{teorema}

\begin{schema}
\Ipotesi{$f$ continua in $[a,b]$; \ $G'=f$ in $[a,b]$.}
\Tesi{$\int_{a}^{b}f=G(b)-G(a)$.}
\end{schema}

\begin{strategia}
La funzione integrale $F(x)=\int_{a}^{x}f$ è una primitiva di $f$ (teorema fondamentale),
come $G$. Per la caratterizzazione delle primitive, $G=F+c$; la costante si determina
calcolando in $x=a$, dove $F$ si annulla, e la tesi si ottiene calcolando in $x=b$.
\end{strategia}

\begin{dimostrazione}
Per il \cref{teo:fondamentale}, $F(x)=\int_{a}^{x}f(t)\,dt$ è una primitiva di $f$ in
$[a,b]$; per il \cref{teo:caratterizzazione-primitive} esiste una costante $c$ tale che
\[
  G(x)=F(x)+c=c+\int_{a}^{x}f(t)\,dt \qquad\perogni x\in[a,b] .
\]
Per $x=a$ si ottiene $G(a)=c+0=c$; sostituendo, $G(x)=G(a)+\int_{a}^{x}f(t)\,dt$, e per
$x=b$ si ha la tesi.
\end{dimostrazione}

Il simbolo $\bigl[G(x)\bigr]_{a}^{b}$ indica la differenza $G(b)-G(a)$; si scrive anche
$G(x)\big|_{a}^{b}$. La formula vale anche se $a>b$, per la \eqref{eq:estremi-invertiti}.

\begin{esempio}[Primi integrali calcolati con la formula fondamentale]
\label{es:formula-fondamentale}
Ecco alcuni integrali che il \cref{cap:integrali-definiti} non permetteva di calcolare, o
permetteva di calcolare soltanto con fatica.
\begin{enumerate}[label=\textup{(\alph*)}]
  \item Il settore di parabola: $G(x)=\frac{x^{3}}3$ è una primitiva di $x^{2}$, e
  \[
    \int_{0}^{b}x^{2}\,dx=\Bigl[\frac{x^{3}}3\Bigr]_{0}^{b}=\frac{b^{3}}3 .
  \]
  Il risultato di Archimede,
  che ci era costato un paragrafo, si ottiene in una riga.
  \item L'arco di sinusoide (\cref{fig:primitive}, a destra):
  $\int_{0}^{\pi}\sen x\,dx=\bigl[-\cos x\bigr]_{0}^{\pi}=-\cos\pi+\cos0=2$.
  \item Poiché $D\operatorname{arctg}x=\frac1{1+x^{2}}$ \eqref{eq:derivata-arctg},
  $\int_{0}^{1}\frac{dx}{1+x^{2}}=\operatorname{arctg}1-\operatorname{arctg}0=\frac\pi4$: è il
  valore annunciato negli esercizi del \cref{cap:integrali-definiti}, e quindi anche
  $\sum_{k=1}^{n}\frac n{n^{2}+k^{2}}\to\frac\pi4$ (\cref{es:limiti-riemann}).
  \item Poiché $D\,2^{x}=2^{x}\log2$ \eqref{eq:derivata-ax}, una primitiva di $2^{x}$ è
  $\frac{2^{x}}{\log2}$, e $\int_{0}^{1}2^{x}\,dx=\frac{2-1}{\log2}=\frac1{\log2}$, come avevamo
  trovato con le progressioni geometriche (\cref{es:integrale-esponenziale}).
\end{enumerate}
\end{esempio}

\begin{osservazione}[La formula fondamentale senza il teorema fondamentale]
\label{oss:formula-fondamentale-integrabili}
La formula \eqref{eq:formula-fondamentale} vale anche se $f$ è soltanto integrabile in
$[a,b]$, purché abbia una primitiva $G$ in $[a,b]$. Infatti, per ogni partizione
$P=\set{x_{0},\dots,x_{n}}$, il teorema di Lagrange applicato a $G$ in ciascun intervallo
$[x_{k-1},x_{k}]$ fornisce punti $\xi_{k}$ con $G(x_{k})-G(x_{k-1})=f(\xi_{k})\,\Delta x_{k}$, e
sommando (la somma a primo membro è telescopica)
\[
  G(b)-G(a)=\sum_{k=1}^{n}f(\xi_{k})\,\Delta x_{k} ,
\]
che è una somma di Riemann, compresa tra $s(P)$ e $S(P)$. Il numero $G(b)-G(a)$ è dunque un
elemento separatore tra le somme inferiori e superiori di $f$, e poiché $f$ è
integrabile coincide con l'integrale. Questa dimostrazione mostra anche perché la
formula funziona: in una qualsiasi partizione, l'incremento totale di $G$ è la somma
degli incrementi sui singoli intervalli, ciascuno uguale a «derivata per ampiezza».
\end{osservazione}

\begin{idea}[Spazio percorso e velocità]
Se $s(t)$ è la posizione di un velivolo lungo la pista e la velocità $v(t)=s'(t)$ è
continua, la formula fondamentale dà
\begin{equation}\label{eq:spazio-velocita}
  s(t_{1})-s(t_{0})=\int_{t_{0}}^{t_{1}}v(t)\,dt :
\end{equation}
lo spazio percorso è l'area sotto il grafico della velocità. È la giustificazione
promessa nel \cref{cap:integrali-definiti} per la corsa di decollo e per il variometro
dell'aliante: la quota $h(t)$ ha derivata $v_{z}(t)$, e la variazione di quota è
$\int v_{z}\,dt$. Allo stesso modo, il teorema della media per gli integrali diventa il
teorema di Lagrange per la primitiva: $\frac1{b-a}\int_{a}^{b}f=\frac{G(b)-G(a)}{b-a}=G'(\xi)=f(\xi)$
per un $\xi\in(a,b)$, come avevamo anticipato.
\end{idea}

\subsection{Approfondimento: la formula fondamentale per le funzioni regolari a tratti}
\label{subsec:regolari-tratti}

Nelle applicazioni si incontrano spesso funzioni continue che sono derivabili tranne che in
un numero finito di punti, dove il grafico ha un punto angoloso: $\abs x$, la funzione
integrale $\abs x-1$ del \cref{sec:teorema-fondamentale}, o, come vedremo, il momento
flettente di un'ala che porta un motore. Seguendo \cite{giusti}, estendiamo a queste
funzioni la formula fondamentale.

\begin{definizione}[Funzione regolare a tratti]\label{def:regolare-tratti}
\index{funzione!regolare a tratti}
Una funzione $G:[a,b]\to\R$ si dice \emph{regolare a tratti} se esistono punti
$a=z_{0}<z_{1}<\dots<z_{r}=b$ tali che, in ciascuno degli intervalli aperti
$(z_{k-1},z_{k})$, $G$ è derivabile con derivata $G'$ continua e limitata.
\end{definizione}

Nei punti $z_{k}$ la derivata può non esistere. Assegnandole in quei punti valori
arbitrari si ottiene però una funzione limitata e continua tranne che in un numero finito
di punti, quindi integrabile (\cref{es:continua-tratti}), il cui integrale non dipende dai
valori scelti (\cref{es:valore-in-un-punto}): è questo il numero che indichiamo con
$\int_{a}^{b}G'(x)\,dx$.

\begin{teorema}[Formula fondamentale per le funzioni regolari a tratti]
\label{teo:formula-regolari-tratti}
Se $G$ è continua in $[a,b]$ e regolare a tratti, allora
\begin{equation}\label{eq:formula-regolari-tratti}
  \int_{a}^{b}G'(x)\,dx=G(b)-G(a) .
\end{equation}
\end{teorema}

\begin{schema}
\Ipotesi{$G$ continua in $[a,b]$; \ $G$ derivabile, con $G'$ continua e limitata, in
ciascun intervallo $(z_{k-1},z_{k})$.}
\Tesi{$\int_{a}^{b}G'(x)\,dx=G(b)-G(a)$.}
\end{schema}

\begin{strategia}
Per l'additività basta ragionare su un intervallo $[z_{k-1},z_{k}]$, negli estremi del
quale $G$ può non essere derivabile. Ce ne teniamo lontani: in un intervallo $[t,s]$ più
piccolo vale la formula ordinaria; poi facciamo tendere $t$ e $s$ agli estremi, usando
la continuità della funzione integrale e quella di $G$, che è dunque essenziale.
\end{strategia}

\begin{dimostrazione}
Fissato $k$, poniamo $\alpha=z_{k-1}$, $\beta=z_{k}$; sia $g$ una funzione uguale a $G'$
in $(\alpha,\beta)$, con valori arbitrari in $\alpha$ e $\beta$, e sia
$\Phi(x)=\int_{\alpha}^{x}g(t)\,dt$ per $x\in[\alpha,\beta]$. Se $\alpha<t<s<\beta$, la
funzione $G$ è derivabile in $[t,s]$ con derivata continua, e per l'additività e il
\cref{teo:formula-fondamentale}
\[
  \Phi(s)-\Phi(t)=\int_{t}^{s}G'(x)\,dx=G(s)-G(t) .
\]
La funzione $\Phi$ è continua in $[\alpha,\beta]$ (\cref{prop:integrale-lipschitz}), e
così pure $G$, per ipotesi. Facendo tendere $t$ ad $\alpha$ da destra si ottiene
$\Phi(s)=G(s)-G(\alpha)$ per ogni $s\in(\alpha,\beta)$, perché $\Phi(\alpha)=0$; facendo
poi tendere $s$ a $\beta$ da sinistra, $\Phi(\beta)=G(\beta)-G(\alpha)$, cioè
$\int_{z_{k-1}}^{z_{k}}G'=G(z_{k})-G(z_{k-1})$. Sommando su $k$ e usando l'additività
(\cref{teo:additivita}), il secondo membro è una somma telescopica che vale $G(b)-G(a)$.
\end{dimostrazione}

\begin{osservazione}[La continuità di $G$ è essenziale]\label{oss:regolari-tratti-salti}
La funzione segno $G$, che vale $-1$ per $x<0$, $0$ in $0$ e $1$ per $x>0$, è regolare a tratti
in $[-1,1]$, con derivata nulla per $x\neq0$: dunque $\int_{-1}^{1}G'=0$, mentre
$G(1)-G(-1)=2$. La differenza è il salto in $0$: la derivata «non vede» i salti. In
generale, se $G$ è regolare a tratti, continua in $a$ e in $b$, e nei punti interni
$z_{k}$ ha limiti destro e sinistro, la stessa dimostrazione (dove ora, per $t\to z_{k-1}^{+}$
e $s\to z_{k}^{-}$, compaiono i limiti laterali di $G$) dà
\[
  \int_{a}^{b}G'(x)\,dx=G(b)-G(a)-\sum_{k=1}^{r-1}\sigma_{k}, \qquad
  \sigma_{k}=\lim_{x\to z_{k}^{+}}G(x)-\lim_{x\to z_{k}^{-}}G(x) .
\]
\end{osservazione}

\begin{esempio}[Una funzione integranda con un salto]\label{es:integranda-salto}
Sia $f(x)=-1$ per $x<0$ e $f(x)=x$ per $x\ge0$. La funzione $f$ non ha primitive in
$[-1,2]$ (se fosse $H'=f$, la derivata di $H(x)+\frac x2$ non si annullerebbe mai e
cambierebbe segno, contro il \cref{lem:darboux}), ma fuori dall'origine è la derivata di
\[
  G(x)=\begin{cases} -x & \text{se } x<0,\\[2pt] \dfrac{x^{2}}2+c & \text{se } x\ge0, \end{cases}
\]
qualunque sia la costante $c$. Con $c=0$ la funzione $G$ è continua e regolare a tratti, e
per il \cref{teo:formula-regolari-tratti} $\int_{-1}^{2}f=G(2)-G(-1)=2-1=1$, come conferma
il calcolo spezzato in $0$. Con $c\neq0$, invece, $G$ ha in $0$ il salto $\sigma_{1}=c$, e
$G(2)-G(-1)=1+c$ è sbagliato proprio di $c$: le costanti dei diversi tratti vanno scelte
in modo da «incollare» i pezzi di $G$ con continuità.
\end{esempio}

\begin{notastorica}[Derivate non integrabili: Volterra e Lebesgue]
Nell'\cref{oss:formula-fondamentale-integrabili} e nel \cref{teo:formula-regolari-tratti}
abbiamo supposto la derivata integrabile. L'ipotesi sembra superflua quando la derivata è
limitata; invece nel 1881 Vito Volterra, allora studente di Ulisse Dini alla Scuola
Normale di Pisa, costruì una funzione derivabile in ogni punto, con derivata limitata ma
non integrabile secondo Riemann, perché discontinua in un insieme che non ha misura
nulla (nel senso della nota storica del \S\ref{sec:integrabilita-continue}): per essa il
primo membro della formula fondamentale non esiste. La questione fu risolta da Henri
Lebesgue: con il suo integrale, introdotto nella tesi del 1902, ogni derivata limitata è
integrabile e la formula $\int_{a}^{b}G'=G(b)-G(a)$ vale sempre. Non a caso le sue lezioni
del 1904 si intitolano \emph{Leçons sur l'intégration et la recherche des fonctions
primitives}: ritrovare una funzione dalla sua derivata è il problema di questo capitolo.
\end{notastorica}
